\worksheettitle
  {Домашняя работа}
  {\modulemark{M01-H} Цифры, числа и разряды}

\studentnameblock

\begin{notebox}[Как выполнять]
  Определяйте разряды справа налево. В ответе различайте саму цифру, её разряд
  и значение. Если это пока трудно, сначала выполните карточку M01-R.
\end{notebox}

\begin{practicebox}[1. Назовите цифру разряда]
  \raggedright
  Сотни в числе \(51\,684\): \answerblank[13mm];\quad
  единицы тысяч в числе \(73\,205\): \answerblank[13mm];\quad
  десятки в числе \(40\,918\): \answerblank[13mm].
\end{practicebox}

\begin{practicebox}[2. Назовите разряд и значение цифры]
  \begin{center}
  \renewcommand{\arraystretch}{1.5}
  \begin{tabular}{c|c|c|c}
    \toprule
    число & цифра & разряд & значение цифры \\
    \midrule
    \(5\,832\) & \(8\) & \answerblank[27mm] & \answerblank[22mm] \\
    \(46\,075\) & \(6\) & \answerblank[27mm] & \answerblank[22mm] \\
    \(90\,204\) & \(2\) & \answerblank[27mm] & \answerblank[22mm] \\
    \bottomrule
  \end{tabular}
  \end{center}
\end{practicebox}

\begin{practicebox}[3. Одинаковые цифры \textemdash{} разные значения]
  \raggedright
  В числе \(72\,727\) цифра \(7\) встречается три раза. Назовите значения
  всех трёх цифр \(7\), двигаясь слева направо:
  \[
    \answerblank[25mm],\qquad \answerblank[25mm],\qquad
    \answerblank[25mm].
  \]
  Почему значения одинаковых цифр различаются? \answerblank[70mm].
\end{practicebox}

\begin{warningbox}[4. Найдите и исправьте ошибку]
  \raggedright
  Ученик сказал: «В числе \(36\,241\) цифра \(6\) имеет значение \(600\)».

  В каком разряде стоит цифра \(6\)? \answerblank[30mm].\par
  Каково её значение? \answerblank[22mm].
\end{warningbox}
